\chapter{Examples of area, critical sets and pushforward densities}\label{sec:examples-area-critical}
The area and coarea formulas in \cite{Sim18}, Chapter 2, Sections 3 and 7, describe different measures even for elementary maps. These original calculations develop their metric factors, critical sets, exceptional fibers and multiplicities. The fat Cantor construction is the one already given in Definition \ref{def:fat-cantor-set}; no new existence theorem about critical sets is assumed.

\begin{example}[A nonorthogonal linear parametrization has a Gram determinant]\label{ex:area-nonorthogonal-linear-map}
Let $F(s,t)=(s,2t,s+t)$ on $(0,1)^2$. Then
\[
 (DF)^*DF=\begin{pmatrix}2&1\\1&5\end{pmatrix},
 \qquad J_2DF=3,
 \qquad \mathcal H^2(F((0,1)^2))=3.
\]
Projection of the image to the first two coordinates has area two. Multiplying the lengths of the two tangent columns would incorrectly give $\sqrt{10}$ for the surface Jacobian.
\end{example}
\begin{proof}
The columns are $(1,0,1)$ and $(0,2,1)$, with squared lengths two and five and inner product one. Their Gram determinant is $10-1=9$. The map is injective, so Lemma \ref{lem:area-positive-metric-factor} gives image area three. The first two coordinates form the rectangle $(0,1)\times(0,2)$, of area two. The product of column lengths is the area of the rectangle they would span if orthogonal; the actual parallelogram includes their nonzero angle through the Gram determinant.
\end{proof}

\begin{example}[A circle traversed several times counts every sheet]\label{ex:area-multiple-circle-traversal}
For an integer $q\geq1$, let $\gamma_q(t)=(\cos(qt),\sin(qt))$, $0\leq t<2\pi$. Then every point of the unit circle has exactly $q$ preimages, and
\[
 \int_0^{2\pi}|\gamma_q'(t)|\,dt=2\pi q
 =q\,\mathcal H^1(S^1).
\]
The image alone has length $2\pi$, independently of $q$.
\end{example}
\begin{proof}
The speed is the constant $q$. Writing a point of $S^1$ by its angle modulo $2\pi$, the equation $qt=\theta$ modulo $2\pi$ has precisely $q$ solutions in the half-open parameter interval. Apply Theorem \ref{thm:euclidean-area-local-proof} to the globally smooth map and that measurable interval. For $q=1$, the formula yields $\mathcal H^1(S^1)=2\pi$ directly; the general case gives the displayed multiplicity factor. Using a half-open interval prevents an extra endpoint preimage at $(1,0)$, although a single target point would not change the length integral.
\end{proof}

\begin{example}[A cusp has a harmless isolated critical parameter]\label{ex:area-real-cusp-length}
The curve $\gamma(t)=(t^2,t^3)$, $-1\leq t\leq1$, is injective and has derivative zero only at $t=0$. Its length is
\[
 \mathcal H^1(\gamma([-1,1]))
 =\int_{-1}^1|t|\sqrt{4+9t^2}\,dt
 =\frac{2}{27}\bigl(13^{3/2}-8\bigr).
\]
Thus failure of immersion at one parameter does not invalidate the area identity.
\end{example}
\begin{proof}
Equality of the second coordinates $t^3$ implies equality of the parameters, proving injectivity. The derivative is $(2t,3t^2)$, with the displayed norm. The locally Lipschitz area formula applies on the closed parameter interval by regarding it as a measurable subset of an open interval. Evenness gives twice the integral on $[0,1]$. Since the derivative of $(4+9t^2)^{3/2}$ is $27t\sqrt{4+9t^2}$, integration gives the stated value. The critical parameter and its image both have zero one-dimensional measure; the geometric cusp at the image is not being declared a smooth point.
\end{proof}

\begin{example}[A strictly increasing Lipschitz map can collapse a positive-measure set]\label{ex:area-fat-cantor-critical-map}
Let $E\subset[0,1]$ be the compact nowhere-dense set of measure $1/2$ in Proposition \ref{prop:fat-cantor-length-topology}. Define
\[
 f(x)=\int_0^x\chi_{[0,1]\setminus E}(s)\,ds,
 \qquad 0\leq x\leq1.
\]
Then $f$ is one-Lipschitz and strictly increasing, maps $[0,1]$ onto $[0,1/2]$, and satisfies $f'=0$ almost everywhere on $E$ and $f'=1$ almost everywhere on its complement. Nevertheless $\mathcal L^1(f(E))=0$. Moreover
\[
 f_\#(\mathcal L^1\llcorner[0,1])
 =\mathcal L^1\llcorner[0,1/2]+\nu,
 \qquad \nu=f_\#(\mathcal L^1\llcorner E),
\]
where $\nu$ is singular continuous and has mass $1/2$.
\end{example}
\begin{proof}
The integral gives absolute continuity, the derivative values almost everywhere, and the one-Lipschitz bound. Every interval of positive length in $[0,1]$ meets the open complement of $E$ in a nonempty interval, because $E$ is nowhere dense. Its integral increment is therefore positive, proving strict increase. The endpoint value is $f(1)=1-|E|=1/2$.

Extend $f$ constantly outside $[0,1]$ to a Lipschitz map. Corollary \ref{cor:area-null-image-criterion} applied to $E$ gives $|f(E)|=0$. Since $f$ is injective, every singleton has at most one preimage; hence $\nu$ has no atoms. Its support lies in the compact null set $f(E)$, and its total mass is $|E|=1/2$, proving singular continuity.

On the complement of $E$, the derivative equals one almost everywhere. The weighted area formula, Proposition \ref{prop:weighted-area}, applied to indicators of target Borel sets and using injectivity, gives
\[
 f_\#(\mathcal L^1\llcorner([0,1]\setminus E))
 =\mathcal L^1\llcorner f([0,1]\setminus E)
 =\mathcal L^1\llcorner[0,1/2].
\]
The last equality ignores only the null set $f(E)$. Adding the contribution from $E$ proves the decomposition. This example produces the singular measure by a Lipschitz map of a restricted positive-volume set; it does not require the usual Cantor function to be Lipschitz.
\end{proof}

\begin{example}[The graph retains the area that a scalar image loses]\label{ex:area-fat-cantor-graph}
For $f$ and $E$ in Example \ref{ex:area-fat-cantor-critical-map}, the graph map $G(x)=(x,f(x))$ is injective and
\[
 \mathcal H^1(G([0,1]))=\frac{1+\sqrt2}{2},
 \qquad \mathcal H^1(G(E))=\frac12,
 \qquad \mathcal L^1(f(E))=0.
\]
A vanishing scalar derivative on $E$ does not make the graph derivative vanish there.
\end{example}
\begin{proof}
The first coordinate makes $G$ injective, and its speed is $\sqrt{1+|f'|^2}$ almost everywhere. This equals one on $E$ and $\sqrt2$ on its complement. The area formula on $[0,1]$ gives $|E|+\sqrt2(1-|E|)=(1+\sqrt2)/2$. Applying it only to $E$ gives length $1/2$. The scalar image assertion was established in the preceding example. In the graph, the derivative has the nonzero first component one even where the second component vanishes.
\end{proof}

\begin{example}[Graph regularization can have a linear Jacobian error]\label{ex:area-linear-regularization-rate}
For an integer $m\geq1$, let $f:\mathbb R^m\to\mathbb R^m$ be $f(x)=(x_1,\ldots,x_{m-1},0)$, with the initial list empty if $m=1$. For $0<\varepsilon\leq1$ and $F_\varepsilon(x)=(f(x),\varepsilon x)$,
\[
 J_mDf=0,
 \qquad J_mDF_\varepsilon
 =\varepsilon(1+\varepsilon^2)^{(m-1)/2}.
\]
In particular this Jacobian is of order $\varepsilon$, and cannot be bounded by $C\varepsilon^2$ with $C$ independent of small $\varepsilon$.
\end{example}
\begin{proof}
The Gram matrix $Df^*Df$ is diagonal with $m-1$ entries equal to one and the last equal to zero. Lemma \ref{lem:area-graph-regularization} adds $\varepsilon^2$ to every entry. The square root of the product gives the displayed expression. Its ratio to $\varepsilon$ tends to one, while its ratio to $\varepsilon^2$ tends to infinity. The linear upper bound used to discard critical images is sufficient, including in this worst-rank example.
\end{proof}

\begin{example}[A Lipschitz interval map has an infinite exceptional fiber]\label{ex:area-infinite-exceptional-fiber}
On $[0,1]$ define $f(0)=0$ and $f(x)=x^2\sin(\pi/x)$ for $x>0$. This map is Lipschitz, and
\[
 f^{-1}(0)=\{0\}\cup\{1/j:j=1,2,\ldots\}.
\]
Thus its multiplicity is infinite at zero, but is integrable with respect to target Lebesgue measure and is finite almost everywhere.
\end{example}
\begin{proof}
For $x>0$,
\[
 f'(x)=2x\sin(\pi/x)-\pi\cos(\pi/x),
 \qquad |f'(x)|\leq2+\pi.
\]
The mean-value theorem gives this Lipschitz constant between positive arguments; continuity at zero passes the bound to an endpoint argument. The zero set follows from the zeros of sine. Extend the map to a Lipschitz function on the line, and apply Proposition \ref{prop:area-integrable-multiplicity} to $[0,1]$. It gives integral at most $2+\pi$ and finite multiplicity almost everywhere. Infinite multiplicity at one null target level is compatible with that conclusion. No boundedness of the multiplicity at every point is claimed by the area formula.
\end{proof}

\begin{example}[Rank loss can give infinite fibers on a null image]\label{ex:area-projection-infinite-fibers}
For $f:(0,1)^2\to\mathbb R^2$, $f(x,y)=(x,0)$, the two-dimensional area Jacobian is zero everywhere. Every point of $(0,1)\times\{0\}$ has an uncountable fiber, yet
\[
 \int_{\mathbb R^2}\#f^{-1}(z)\,d\mathcal L^2(z)=0.
\]
The integrand is infinity on a null line segment and zero elsewhere.
\end{example}
\begin{proof}
The derivative matrix has rank one, so its determinant and area Jacobian vanish. At an image point $(x,0)$ the fiber is $\{x\}\times(0,1)$, whose counting measure is infinite. The image line segment has zero two-dimensional measure by diameter covers. A nonnegative extended-valued function supported on a null set has integral zero, regardless of infinite values there. This computes both sides of the area formula and separates target measure from the cardinality of an individual fiber.
\end{proof}

\begin{example}[A branched power map counts degree in its area]\label{ex:area-complex-power-map}
For integer $q\geq1$ and $R>0$, let $f(z)=z^q$ on the complex disk $D_R$. Its real two-dimensional Jacobian is
\[
 J_2Df(z)=q^2|z|^{2q-2},
 \qquad \int_{D_R}J_2Df\,dA=\pi qR^{2q}.
\]
The image disk has area $\pi R^{2q}$; its nonzero points have $q$ preimages in $D_R$.
\end{example}
\begin{proof}
At $z\ne0$, the complex derivative is $qz^{q-1}$, which acts on the real plane as a rotation followed by dilation by its modulus. Its real area Jacobian is the square of that modulus. The same formula holds at zero, with value one if $q=1$ and zero if $q>1$. Polar integration gives
\[
 2\pi q^2\int_0^Rr^{2q-1}\,dr=\pi qR^{2q}.
\]
Every nonzero image value has exactly the $q$ distinct complex roots, all of modulus less than $R$; the value zero has only one preimage. The area formula therefore integrates multiplicity $q$ on the image disk except for a null target point. This is a branched parametrization of a plane; it should not be confused with an injectively parametrized singular complex curve in higher codimension.
\end{proof}

\begin{example}[A product map has a logarithmically unbounded density]\label{ex:coarea-product-log-density}
For $f(x,y)=xy$ on $Q=(-1,1)^2$,
\[
 f_\#(\mathcal L^2\llcorner Q)
 =-2\log|t|\,\chi_{\{0<|t|<1\}}\mathcal L^1.
\]
The density is defined by the displayed expression for $0<|t|<1$ and by zero elsewhere; its value at zero may be changed freely. It is unbounded near zero but integrable, with total mass four. There is no atom at zero, although that fiber is the union of two crossing segments.
\end{example}
\begin{proof}
For a nonnegative Borel test function $\phi$, ordinary one-dimensional substitution at each fixed $x\ne0$ gives
\[
 \int_{-1}^1\int_{-1}^1\phi(xy)\,dy\,dx
 =\int_{-1}^1\frac1{|x|}\int_{-|x|}^{|x|}\phi(t)\,dt\,dx.
\]
Tonelli makes the density at $0<|t|<1$ equal to $2\int_{|t|}^1dx/x=-2\log|t|$. Its integral is four since $\int_0^1-\log t\,dt=1$. The zero fiber consists of the two coordinate segments and has zero area, so it carries no pushforward atom. Away from $x=0$, parametrize a fiber by $(x,t/x)$. Its speed divided by $|Df|$ equals
\[
 \frac{\sqrt{1+t^2/x^4}}{\sqrt{x^2+t^2/x^2}}=\frac1{|x|},
\]
which gives the same inverse-Jacobian density by Proposition \ref{prop:coarea-inverse-jacobian}. The crossing at zero affects one exceptional level only.
\end{proof}

\begin{example}[Radial powers give integrable power singularities]\label{ex:coarea-radial-power-density}
Let $m\geq2$, $R>0$, $\alpha\geq1$, and $f(x)=|x|^\alpha$ on $B_R(0)\subset\mathbb R^m$. Then
\[
 f_\#(\mathcal L^m\llcorner B_R)
 =\frac{m\omega_m}{\alpha}\,t^{m/\alpha-1}
   \chi_{(0,R^\alpha)}(t)\mathcal L^1.
\]
When $\alpha>m$ this density diverges at zero, but its integral is $\omega_mR^m$ and it has no atom there.
\end{example}
\begin{proof}
The map is locally Lipschitz for $\alpha\geq1$. At a positive level $t$, its fiber is the sphere of radius $r=t^{1/\alpha}$, with area $m\omega_mr^{m-1}$, and its gradient has norm $\alpha r^{\alpha-1}$. Their ratio is $(m\omega_m/\alpha)t^{m/\alpha-1}$, as in Proposition \ref{prop:coarea-inverse-jacobian}. Equivalently the sublevel volume is $\omega_m t^{m/\alpha}$, giving the same derivative. Its value tends to zero at the origin, which also shows there is no atom. The exponent $m/\alpha-1$ is always greater than $-1$, so the singularity remains integrable; integration up to $R^\alpha$ yields the original ball volume.
\end{proof}

\begin{example}[Distance slices of an offset affine plane]\label{ex:coarea-offset-plane-distance}
Let $m\geq2$, $d>0$, $P=\mathbb R^m\times\{d\}\subset\mathbb R^{m+1}$, and $\rho(x)=|x|$. For $r>d$,
\[
 \mathcal H^m(P\cap B_r)=\omega_m(r^2-d^2)^{m/2},
 \qquad
 \frac{d}{dr}\mathcal H^m(P\cap B_r)
 =m\omega_mr(r^2-d^2)^{(m-2)/2}.
\]
The derivative equals the inverse-Jacobian fiber integral. The unique closest point is critical and contributes no atom.
\end{example}
\begin{proof}
In the coordinates $x=(u,d)$ of the plane, the condition $|x|<r$ becomes $|u|<\sqrt{r^2-d^2}$. This gives the volume and its derivative. On the fiber, the tangential gradient norm is $|u|/r=\sqrt{r^2-d^2}/r$, while its $(m-1)$-measure is $m\omega_m(r^2-d^2)^{(m-1)/2}$. Dividing gives the displayed derivative. The only critical point is $u=0$, a set of zero $m$-measure. The volume tends to zero as $r\downarrow d$ and has no jump there. This contrasts with a whole sphere centered at the origin, where every point lies in one critical distance level.
\end{proof}
